% =====================================================================
%  CONTEXTUALIZACIÓN
%  Los objetivos se transcriben LITERAL del anteproyecto §4.1 y §4.2: son los
%  aprobados y no se reescriben.
% =====================================================================
\chapter{Contextualización}

\section{Planteamiento del problema}

Orientarse en un edificio desconocido exige construir y actualizar una representación
mental del espacio, un proceso que se complica cuando el entorno tiene visibilidad
limitada, estructuras repetidas y varios puntos de decisión. Las edificaciones de
varios pisos añaden una dificultad propia, debido a que el cambio de piso rompe la
continuidad espacial y obliga a integrar información vertical en el recorrido.

Aunque la robótica móvil colaborativa ha avanzado en exploración, logística y
cobertura distribuida, la mayoría de esos desarrollos se evaluó en entornos de un solo
piso, donde la segmentación vertical no es el problema central. Los casos de limpieza
doméstica y de almacén que documenta \textcite{nair2024collaborative} son
representativos de ese escenario.

Resolver el cambio de piso con un solo robot, por su parte, resulta costoso, ya que
exige mecanismos de locomoción capaces de subir escaleras o intervenir el sistema de
control del ascensor. Queda abierta, entonces, la manera de asistir a un usuario que
cambia de piso sin encarecer la plataforma ni modificar el edificio.

\section{Pregunta problema}

\begin{quote}
¿Cómo asistir la orientación de usuarios dentro de entornos interiores con
múltiples pisos utilizando sistemas robóticos?
\end{quote}

\section{Objetivos}

\subsection{Objetivo general}

Implementar un sistema colaborativo de dos robots móviles terrestres para la guía
coordinada y eficiente de asistencia a visitantes en entornos interiores con
múltiples pisos, mediante comunicación inter-robot, asignación dinámica de tareas
y una interfaz móvil de interacción usuario-robot.

\subsection{Objetivos específicos}

\begin{enumerate}[leftmargin=*,label=\textbf{\arabic*.}]
  \item Modelar la arquitectura funcional del sistema colaborativo de dos robots
    móviles, definiendo requerimientos, escenarios de operación con múltiples
    pisos, estrategia de asignación de tareas y esquema de comunicación
    inter-robot.
  \item Desarrollar una plataforma robótica móvil basada en dos vehículos tipo
    DonkeyCar, integrando los módulos de locomoción, sensado, procesamiento y
    comunicación necesarios para la operación colaborativa en entornos
    interiores.
  \item Programar una interfaz móvil para interacción humano-robot (HRI) basada
    en la selección de localizaciones de interés de origen y destino.
  \item Evaluar el desempeño del sistema mediante pruebas experimentales en un
    entorno interior controlado, utilizando métricas como tiempo de respuesta,
    tiempo de asignación de robot, tasa de éxito en la entrega de asistencia y
    continuidad del servicio entre niveles.
\end{enumerate}

